\section{Three gates, implemented and measured}
\label{sec:gates}

The checks in Section~\ref{sec:caught} were done by hand, and a discipline that
depends on someone deciding to look again is not a discipline. We implemented
all three and ran each on the defective code and on the first correction.
Constraint-based data validation checks declared expectations about data at
scale \citep{schelter2018deequ}; these gates instead check the pipeline's reading
of a vocabulary, which is where the defect lived. The second error gives a test
the first did not: the gates were built after the first error was known, but the
second was unknown when they were built, so how they behave on it is closer to a
prospective measurement. Table~\ref{tab:gates} gives the result.

\begin{table}[t]
\centering
\small
\caption{Each gate on the defective code, and on the first correction, which
still contained the second error of Section~\ref{sec:second} (42 objects).}
\label{tab:gates}
\begin{tabular}{>{\raggedright\arraybackslash}p{2.5cm}>{\raggedright\arraybackslash}p{3.9cm}>{\raggedright\arraybackslash}p{5.4cm}}
\toprule
Gate & Defective code & First correction, against the second error \\
\midrule
Total-function check on the status column &
\textbf{Fires}: 9 unclassified codes over 1{,}099 rows (\texttt{ATT},
\texttt{C}, \texttt{DK}, \texttt{E}, \texttt{EVA DP}, \texttt{GRP}, \texttt{N},
\texttt{REL}, \texttt{TFR}) &
\textbf{Silent}: 0 unclassified codes. \texttt{E} and \texttt{C} are classified,
wrongly \\
\addlinespace
Residue decomposition &
\textbf{Shows it}: 20 groups, top 5 hold 729 of 932; 500 pair docking or
attachment with impact or landing &
\textbf{Shows it, unread}: 44 of 261 in groups pairing \texttt{E} or \texttt{C}
with a CelesTrak object in orbit \\
\addlinespace
Source inventory reconciliation &
\textbf{Fires}: \texttt{satcat100k} listed, never fetched &
\textbf{Silent}: scoped to the four main catalogues; the event catalogue is on
the same index \\
\bottomrule
\end{tabular}
\end{table}

\paragraph{The total-function check tests completeness, not meaning.}
Requiring every observed status to be assigned to a named category attacks the
default that did the damage in Section~\ref{sec:mechanism}, and on the defective
code it lists every code implicated in the error. It cannot tell a right
assignment from a wrong one. The first correction assigned \texttt{E} and
\texttt{C} explicitly, to the wrong category, and the check passed. It also has a
scoping cost: applied to every controlled column rather than the one the
classifier consumes, it fires identically on both versions on 70,292 rows of
CelesTrak's orbit type, 70,130 of object type and 1,292 of data status, columns
the pipeline treats as opaque strings. The scope need not be chosen by hand. The
names both gate paths import from one module are exactly the definitions the
gate cannot test (Section~\ref{sec:mechanism}); at the published commit the
vocabularies among them are all classifications of the GCAT status column, which
is where the check separates cleanly.

\paragraph{Residue decomposition is an auditing aid, not a verdict.}
Grouping the unexplained disagreements by the two registers' fields made the
first error legible, since a status meaning ``joined another object'' should not
systematically meet a landing in the other register. It made the second error
legible too: on the corrected code, 30 explosions and 14 collisions sit against
CelesTrak objects in orbit. We read that table after the correction and saw
confirmation. The decomposition has no rejection criterion of its own, and the
reader supplied the wrong one.

\paragraph{Inventory reconciliation is exactly as good as its list.}
Comparing the files the source lists against the files the harvest loaded found
the missing fourth catalogue at once, with no false positive. On the second error
it was silent, because we had scoped it to the section of the index headed as
the main catalogues, while the index also lists the event catalogue, and eleven
object catalogues in all. Widening the list would have caught it and would also
flag files the pipeline has no reason to read. Running this gate also corrected
our own build report, which said the missing file recovered 280 objects; it
recovers 278, forced by the 900 and 622 that both reproduce exactly. A pass
undertaken to correct three published numbers introduced a fourth error.

\paragraph{A check that looked independent and was not.}
The SHACL validation returned 3{,}564 results, exactly the sum of six of the
seven counts. At the time this read as corroboration. It is an arithmetic
identity: the shapes select the nodes the counts count, from the graph the counts
were computed over, so when the counts were wrong the total was wrong by the same
amount and still matched. A suite of such steps produces a great deal of green
while measuring one thing repeatedly.
